\section{From supersaturation to a finite graph}\label{sec:distributions}

A \emph{unit} is an ordered pair of raw vertices with no hole between its endpoints.
The endpoints of a unit may be dependent.
When two units are sampled independently, their four endpoints need not be independent within either unit.
The event of interest is that all four cross pairs have holes, as in \Cref{fig:four-holes}.

\begin{theorem}[Raw supersaturation]\label{thm:raw-supersaturation}
For the construction parameters specified in \Cref{sec:geometry,sec:parameters}, for all sufficiently large \(n\), every probability law \(\sigma\) on units satisfying
\begin{equation}\label{eq:raw-law-caps}
 \sigma_1\le M\mu,\qquad \sigma_2\le M\mu,\qquad
 \sigma\le 2^{DN}\mu^2
\end{equation}
has the following property: two independent units drawn from \(\sigma\) have all four cross holes with probability at least \(2^{-100gN}\).
The inequalities in \Cref{eq:raw-law-caps} are pointwise inequalities of measures on the finite raw spaces.
\end{theorem}

The proof of \Cref{thm:raw-supersaturation} occupies the subsequent sections and is completed in \Cref{sec:rare-overlap}.
We prove its implication for graphs here. The frame estimates and the preparation of an admissible unit law begin in \Cref{sec:frame-laws}.

\subsection{Entropy minimizers and terminal cuts}

We use a finite form of the information-projection argument of
Csisz\'ar~\citep{csiszar1975}. Its support assertion is important here:
the minimizing law must see every point that can occur in a feasible law.
Relative entropy is taken with natural logarithms:
\[
 \mathsf D(\rho\Vert q)=\sum_x\rho(x)\log\frac{\rho(x)}{q(x)}.
\]
Here \(q\) is a strictly positive probability measure and \(0\log0=0\).

\begin{lemma}\label{lem:entropy-support}
Let \(\mathcal P\) be a nonempty compact convex set of probability measures on a finite set, and let \(\rho\) minimize \(\mathsf D(\cdot\Vert q)\) on \(\mathcal P\).
Then \(\rho\) is positive on the union of the supports of measures in \(\mathcal P\).
For every \(\rho'\in\mathcal P\),
\[
 \mathsf D(\rho'\Vert q)-\mathsf D(\rho\Vert q)
 \ge \mathsf D(\rho'\Vert\rho).
\]
If \(\rho'\) is supported on a set \(S\), the right side is at least \(-\log\rho(S)\).
\end{lemma}
\begin{proof}
If \(\rho(x)=0<\rho'(x)\) at a feasible support point, the entropy along
\((1-t)\rho+t\rho'\) has right derivative \(-\infty\) at zero from the newly occupied points.
At all points where \(\rho>0\), its derivative is finite.
This contradicts minimality, proving the support assertion.

We can therefore differentiate toward any feasible \(\rho'\) using only points where \(\rho>0\).
The right derivative at the minimizer is nonnegative, giving
\[
 \sum_x(\rho'(x)-\rho(x))\log\frac{\rho(x)}{q(x)}\ge0,
\]
where the constant derivative term cancels because both measures have total mass one.
The exact identity
\[
 \mathsf D(\rho'\Vert q)-\mathsf D(\rho\Vert q)
 =\mathsf D(\rho'\Vert\rho)
  +\sum_x(\rho'(x)-\rho(x))\log\frac{\rho(x)}{q(x)}
\]
proves the second assertion.
Finally, if \(\rho'\) is supported on \(S\), comparison with the normalized restriction \(\rho(\cdot\mid S)\) gives
\[
 \mathsf D(\rho'\Vert\rho)
 =\mathsf D(\rho'\Vert\rho(\cdot\mid S))-\log\rho(S)
 \ge-\log\rho(S).
\]
For completeness, nonnegativity of relative entropy follows from
\(\log t\le t-1\), applied with \(t=q(x)/\rho'(x)\) and summed over the positive support of \(\rho'\).
\end{proof}

The next lemma is the max-flow/min-cut principle of
Ford and Fulkerson~\citep{fordfulkerson1956} with capacities chosen
to encode the three density bounds. We include the finite residual-cut
argument, which applies to the real capacities used here.

\begin{lemma}\label{lem:terminal-cut}
Let \(R\subseteq\Omega_n^2\).
If no probability law supported on \(R\) satisfies \Cref{eq:raw-law-caps}, there are sets \(S\subseteq\Omega_n\) and \(E_0\subseteq\Omega_n^2\) such that
\[
 \mu(S)<1/M,\qquad \mu^2(E_0)<2^{-DN},
\]
and every pair in \(R\) either has an endpoint in \(S\) or belongs to \(E_0\).
\end{lemma}
\begin{proof}
Construct a directed network with source \(s\), a left copy and a right copy of \(\Omega_n\), and sink \(t\).
Give the edge from \(s\) to left \(x\) capacity \(M\mu(x)\), the edge from right \(y\) to \(t\) capacity \(M\mu(y)\), and the edge from left \(x\) to right \(y\) capacity
\(2^{DN}\mu(x)\mu(y)\) whenever \((x,y)\in R\).
A flow of value one is precisely a law satisfying the three caps.
A flow of larger value could be scaled down to value one, so the maximum value is less than one.

The feasible flows form a compact polytope, hence the flow value has a maximizer \(f\).
In its residual network, include a forward edge whenever its capacity exceeds its flow and a reverse edge whenever its flow is positive.
There is no residual path from \(s\) to \(t\): augmenting by the smallest positive residual capacity on such a path would increase the value.
Let \(Z\) be the vertices reachable from \(s\).
Every forward edge across \(Z\) to its complement is saturated, and every forward edge in the opposite direction has zero flow, since otherwise its reverse residual edge would leave \(Z\).
Summing flow conservation over \(Z\) shows that the capacity of this cut equals the value of \(f\), hence is less than one.

Let \(L_0\) be the raw types whose left copies are outside \(Z\), and let \(R_0\) be those whose right copies are in \(Z\).
The cut-capacity inequality is
\[
 M\mu(L_0)+M\mu(R_0)
 +2^{DN}\mu^2\bigl(R\cap((\Omega_n\setminus L_0)
                         \times(\Omega_n\setminus R_0))\bigr)<1.
\]
Take \(S=L_0\cup R_0\) and take \(E_0\) to be the pair set in the last summand.
Its measure and that of \(S\) have the asserted bounds.
A pair in \(R\) with neither endpoint in \(S\) necessarily belongs to \(E_0\).
\end{proof}

\subsection{Containers and the random sample}

We now adapt the independent-set fingerprint method of
Kleitman and Winston~\citep{kleitmanwinston1982}; see also
\citet[Section~2]{samotij2015}. Here the progress measure is the minimum
relative entropy of a feasible law. The preceding cut lemma handles
the terminal set when no such law remains.

\begin{proposition}\label{prop:raw-to-graph}
Assume the conclusion of \Cref{thm:raw-supersaturation} at a given sufficiently large \(n\).
For \(m=2^{C_0gN}\), the sampled graph defined in \Cref{sec:geometry} satisfies
\[
 \alpha(G)\le2,\qquad \cm(G)<m/100
\]
with probability \(1-\exp(-\Omega(m))\).
The implied positive constant is independent of \(n\).
\end{proposition}
\begin{proof}
\textbf{Entropy-controlled fingerprints.}
Make a finite \emph{conflict graph} whose vertices are raw units; two units are adjacent when they give all four cross holes.
There are no loops, since a repeated unit would require a hole from a raw vertex to itself.
Put \(\varepsilon_N=2^{-100gN}\).

For an independent set \(I\) of raw units, start with all units as a residual set \(R\).
If there is a feasible law satisfying \Cref{eq:raw-law-caps} on \(R\), choose its relative-entropy minimizer with respect to \(\mu^2\).
Such a minimizer exists because the feasible laws form a compact convex polytope and relative entropy is continuous on the finite probability simplex.
Fix deterministic choices throughout.
The average, under this law \(\rho\), of \(\rho(N_R(v))\) is the conflict probability of two independent \(\rho\)-units, so is at least \(\varepsilon_N\).
Choose a unit \(v\in R\) maximizing this neighborhood mass.
If \(v\in I\), record \(v\) and delete its entire neighborhood \(N_R(v)\).
If \(v\notin I\), delete \(v\) alone.
Both operations preserve \(I\subseteq R\).
Each removes at least one unit, so the procedure terminates at a residual set with no feasible law.

When a recording leaves a feasible residual, \Cref{lem:entropy-support} shows that the new minimum entropy exceeds the old one by at least
\(-\log(1-\varepsilon_N)\).
Other deletions cannot decrease the minimum.
Every feasible law has entropy in \([0,DN\log2]\), since its density is at most \(2^{DN}\).
Allowing one final recording that leaves no feasible residual, the fingerprint has length at most
\begin{equation}\label{eq:fingerprint-length}
 L_N=1+\left\lceil
       \frac{DN\log2}{-\log(1-\varepsilon_N)}
      \right\rceil
 =O(1+DN\,2^{100gN}).
\end{equation}

The fingerprint determines the terminal set.
Indeed, after a unit is recorded its residual neighborhood is empty, and remains empty as residual sets shrink.
A feasible law always has a unit with strictly positive neighborhood mass, so a recorded unit is never selected again.
Given the set of recorded units, replay the deterministic procedure, recording a selected unit exactly when it is in that set.
All other selected units are deleted.
This reproduces every step and the terminal set.

By \Cref{eq:raw-space-size}, the logarithm of the number of raw units is \(O(N^2)\).
The number of possible fingerprints, and hence terminal sets, is at most
\((|\Omega_n|^2+1)^{L_N}\).
Its logarithm is
\begin{equation}\label{eq:container-count}
 O(N^3\,2^{100gN})=o(m),
\end{equation}
since \(m=2^{1000gN}\).

\medskip\noindent\textbf{Terminal exceptions in the sample.}
For each terminal set fix sets \(S,E_0\) supplied by \Cref{lem:terminal-cut}.
Let \(k=\lceil m/200\rceil\).
The probability that at least \(k\) sampled positions have raw type in \(S\) is at most
\[
 \binom mk\mu(S)^k
 \le \left(\frac{em}{kM}\right)^k
 \le (200e/M)^k
 =\exp(-\Omega(m)).
\]
A fixed collection of \(k\) disjoint ordered pairs of distinct sampled positions consists of independent \(\mu^2\)-draws.
It lies entirely in \(E_0\) with probability at most \(2^{-kDN}\).
There are at most \(m^{2k}\) such collections, so the probability that one exists is at most
\begin{equation}\label{eq:pair-exception-bound}
 m^{2k}2^{-kDN}
 =2^{(2C_0g-D)Nk}
 =2^{-2C_0gNk}.
\end{equation}
Union over the \(\exp(o(m))\) terminal sets using \Cref{eq:container-count}.
Outside a set of probability \(\exp(-\Omega(m))\), every terminal set has fewer than \(m/200\) exceptional positions and fewer than \(m/200\) disjoint position pairs in its pair exception.
The strict bounds follow from \(k-1<m/200\).

Suppose on this event that \(G\) has a touching matching of size at least \(m/100\).
Orient its edges arbitrarily, and let \(I\) be the set of raw unit types thereby realized.
This is independent in the conflict graph: four cross holes would mean that the two endpoint sets have no edge between them in \(G\).
Different sampled matching edges may have the same raw type; no injectivity of this assignment is needed.
Run the container procedure for this set \(I\).
Every occurrence of a matching-edge type lies in its terminal set, and hence is covered by \(S\) or \(E_0\).
Because the matching uses disjoint sampled positions, fewer than \(m/200\) of its edges can meet the positions in \(S\).
The rest give disjoint ordered position pairs in \(E_0\), of which there are also fewer than \(m/200\).
Their sum is less than \(m/100\), a contradiction.

Finally, \(\alpha(G)\le2\) holds deterministically by \Cref{lem:hole-triangle-free}.
\end{proof}

\subsection{The clique-minor bound}

\Cref{prop:raw-to-graph} supplies the required bound on connected
matchings once raw supersaturation is proved. The next counting
argument converts that bound into the clique-minor estimate used in
the introduction.

The following elementary estimate counts small branch sets, in the
same connected-matching approach used by Cambie~\citep{cambie2021}.
Its deliberately coarse bound holds for every finite graph.

\begin{proposition}\label{prop:matching-minor}
For every finite nonempty graph \(G\) of order \(m\),
\[
 \had(G)\le\frac{m+4\cm(G)+2}{3}.
\]
In particular, if \(\alpha(G)\le2\) and \(\cm(G)<m/100\), then
\[
 \had(G)<\frac{26m}{75}+\frac23<\frac m2\le\chi(G)
 \qquad(m\ge5).
\]
\end{proposition}
\begin{proof}
Consider a complete-minor model with \(b\) branch sets.
Let \(s\) of them be singletons and \(e\) have two vertices.
Each two-vertex branch is an edge.
The singleton vertices form a clique, so pairing them gives \(\lfloor s/2\rfloor\) more edges.
All these edges together form one touching matching: their vertex sets are disjoint, and adjacency of the original branches guarantees every required contact.
Writing \(c=\cm(G)\), we obtain
\[
 c\ge e+\lfloor s/2\rfloor,\qquad s\le2(c-e)+1.
\]
Every remaining branch has at least three vertices, whence
\[
 m\ge s+2e+3(b-s-e)=3b-2s-e.
\]
Combining the inequalities gives
\[
 3b\le m+2s+e\le m+4c-3e+2\le m+4c+2.
\]
Maximize over all complete-minor models.
For the final assertion, every color class has at most two vertices when \(\alpha(G)\le2\), so \(\chi(G)\ge m/2\).
Combine this with \(c<m/100\) and \(26m/75+2/3<m/2\) when \(m\ge5\).
\end{proof}
